\documentclass[11pt]{report}
\input{../../../report.tex}
\title{Algebra Exercises}
\usepackage{titlepageBU}
\usepackage{xparse}
\courseID{CAS MA 741}
\college{College of Arts and Sciences}
\department{Department of Mathematics and Statistics}
\date{Spring 2025}
\usepackage{tikz-cd}
\DeclareMathOperator{\Stab}{Stab}
\DeclareMathOperator{\rank}{rk}
\declaretheoremstyle[headfont=\bfseries\color{blue}, bodyfont=\normalfont, mdframed={linewidth = 1pt}]{x}
\declaretheorem[numberwithin=section, style=x, name=Exercise]{exercise}
\declaretheorem[numbered=no, style=proof, name=Solution]{solution}
\begin{document}
\makereport
\stepcounter{chapter}
\stepcounter{chapter}
\stepcounter{chapter}
\stepcounter{chapter}
\stepcounter{chapter}
\stepcounter{chapter}
\stepcounter{chapter}
\stepcounter{chapter}
\chapter{Homological Algebra}
\section{(Un)necessary categorical preliminaries}
\begin{exercise}
	Show that if $\psi \phi $ is epi, then $\psi  $ is epi. Show that if $\psi \phi $ is mono, then $\phi $ is mono.
\end{exercise}
\begin{proof}
	Let $\psi \circ \phi $ be epi. Let $\alpha \circ \psi =\beta \circ \psi $. Then $\alpha \circ \psi \circ \phi =\beta \circ \psi \circ \phi $. Since $\psi \circ \phi $ is epi, $\alpha =\beta $. Therefore, $\psi $ is epi. Since the statement for monomorphisms is the dual of that for epis, we are done.
\end{proof}
\begin{exercise}
	Let $\phi : A \to B$ be a morphism in an additive category. Prove that $-\phi $ is a monomorphism or epimorphism if and only if $\phi $ is a monomorphism or epimorphism, respectively.
\end{exercise}
\begin{proof}
	I prove the statement for monomorphisms, since epimorphisms are proved equivalently. Let $\phi $ be mono. Since $\mathsf{A} $ is additive, $-\phi $ is mono if and only if $-\phi \circ \zeta =0$ implies $\zeta =0$. Let $-\phi \circ \zeta =0$. It will be shown later that $0\circ \phi =0$ for any morphism $\phi $, without using this result. It follows that
	\begin{align*}
		-\phi \circ \zeta &=0\\
				  &=0\circ \zeta \\
				  &=(\phi -\phi )\circ \zeta \\
				  &=\phi \circ \zeta -\phi \circ \zeta 
	.\end{align*}
	Adding $\phi \circ \zeta $ to both sides, we get $\phi \circ \zeta =0$. Since $\phi $ is mono, $\zeta =0$. The converse is proven the same way.
\end{proof}
\begin{exercise}
	Preadditive categories are categories where each Hom-set has an abelian group structure, and morphism composition is bilinear with respect to the group structures. Prove that a ring is a preadditive category with one element.
\end{exercise}
\begin{proof}
	Ring elements are morphisms. Since $\End_{\mathsf{A} }(*) $ is endowed with an abelian group structure, we get the underlying group for free. Similarly, morphism composition and the existence of an identity element come for free from category axioms, so we can define multiplication as morphism composition. It suffices to show distributivity holds, and indeed, since $\_\circ \phi $ and $\phi \circ \_$ are linear for any fixed morphism $\phi $, we have
	\[
		\left( \psi + \eta  \right) \circ \phi =(\psi \circ \phi )+(\eta \circ \phi )
	\]
	and
	\[
		\phi \circ \left( \psi +\eta  \right) =\left( \phi \circ \psi  \right) +\left( \psi \circ \eta  \right) 
	.\]
	Therefore, the structure of $\End_{\mathsf{A} }(*) $ is precisely that of a ring.
\end{proof}
\begin{exercise}
	Let $\mathsf{A} $ be a (pre)additive category, and let $A\in \Obj(\mathsf{A} ) $. Show that $\End_{\mathsf{A} }(A) $ has a natural ring structure.
\end{exercise}
\begin{proof}
	By the same argument as the last exercise, we get the abelian group structure, associativity of multiplication, and an identity under multiplication for free, and distributivity follows from bilinearity of composition wrt addition.
\end{proof}
\begin{exercise}
	Let $A,B$ be objects of an additive category $\mathsf{A} $, with zero object 0. Since 0 is initial and final, $\Hom_{\mathsf{A} }(A, 0) $ and $\Hom_{\mathsf{A} }(0, B) $ are singletons. Show that the composition
	\[\begin{tikzcd}[row sep = 2.5em, column sep = 2.5em]
		A\arrow[r]&0\arrow[r]&B,
	\end{tikzcd}\]
	denoted $e$ in this problem but $0$ elsewhere in the book, is the additive identity of $\Hom_{\mathsf{A} }(A, B) $.

	Prove that for any morphism $\phi $ in $\mathsf{A} $, $0\circ \phi =\phi \circ 0=0$.
\end{exercise}
\begin{proof}
	We show $e+e=e$, and this implies $e $ is the identity by basic group theory. Let $a : A \to 0$, $b : 0 \to B$ be the unique morphisms in question, and note that $e=ba$. Since these sets are singletons, we must have $a,b$ as the additive identities of their respective Hom-sets, so $a+a=a$ and $b+b=b$. We have
	\[
		e+e=ba+ba=b(a+a)=b+a=e
	.\]
	
	For the proof that $0\circ \phi =\phi \circ 0=0$, we can simply do the same proof in group theory:
	\[
		\phi \circ 0=\phi \circ (0+0)=\phi \circ 0+\phi \circ 0
	.\]
	Cancelling $\phi \circ 0$ on both sides, we have $\phi \circ 0=0$. The same argument applies in the reverse direction.
\end{proof}
\begin{exercise}
	Prove that an object $A$ of an additive category $\mathsf{A} $ is a zero-object if and only if $1_A$ is the zero-morphism $A\to A$.
\end{exercise}
\begin{proof}
	If $A$ is the zero object, then the statement is immediate since $\End_{\mathsf{A} }(A) $ is a singleton, so suppose that $1_A=0$. We will prove $A$ is the zero-object. We immediately have that $\End_{\mathsf{A} }(A) $ is a singleton, since $1_A$ is both the additive and multiplicative identity under the ring structure uncovered in exercise 9.1.4. In particular, $1_A=0$. Let $\phi : A \to B$ for some object $B$. We then have $\phi =\phi 1_A=\phi 0=0$ by exercise 9.1.5. Therefore, every morphism $\phi : A \to B$ is $0$, and the set is a singleton. The other direction is the same; let $\psi : B \to A$, then $\psi =1_A \psi =0\psi =0$.
\end{proof}
\begin{exercise}
	Give a categorical definition for cokernel in the same spirit as the definition of a kernel as the equalizer of a morphism with $0$, and verify that the definition given here agrees with that requirement.
\end{exercise}
\begin{proof}
	Recall the definition given earlier: let $\mathsf{I} $ be the category
	\[\begin{tikzcd}[row sep = 2.5em, column sep = 4em]
		\mathbf{1} \arrow[r, bend left, "\alpha "]\arrow[r, bend right, "\beta "']&\mathbf{2}.
	\end{tikzcd}\]
	Let $\mathscr{F} : \mathsf{I}  \to \mathsf{A} $ be a functor such that $\mathscr{F} (\alpha )=\phi $ and $\mathscr{F} (\beta )=0$. The corresponding diagram in $\mathsf{A} $ is thus
	\[\begin{tikzcd}[row sep = 2.5em, column sep = 4em]
		A_1\arrow[r, bend left, "\phi "]\arrow[r, bend right, "0"']&A_2.
	\end{tikzcd}\]
	We then were able to define $\ker \phi =\varprojlim \mathscr{F} $. This implies that $\coker \phi $ would be the dual of this notion; that is, $\coker \phi =\varinjlim \mathscr{F} $. This turns out to be precisely the case. Let $(C,\eta ,\psi )$ be a co-cone of $\mathscr{F} $. Then
	\[\begin{tikzcd}[row sep = 2.5em, column sep = 2.5em]
		A_1\arrow[dr, "c_1"']\arrow[ddr, bend right, "\eta "']\arrow[rr, "0"]&&A_2\arrow[dl, "c_2"]\arrow[ddl, bend left, "\psi "]\\
										   &\varinjlim \mathscr{F} \arrow[d, dashed]\\
										   &C
	\end{tikzcd}\]
	commutes. It follows that $\eta ,c_1=0$, and moreover for any co-cone $C$ of $\mathscr{F} $, the map $A_1\to C$ must be $0$. Plugging this into the diagram for $\phi $, we have
	\[\begin{tikzcd}[row sep = 2.5em, column sep = 2.5em]
		A_1\arrow[dr, "0"']\arrow[ddr, bend right, "0"']\arrow[rr, "\phi "]&&A_2\arrow[dl, "c_2"]\arrow[ddl, bend left, "\psi "]\\
										   &\varinjlim \mathscr{F} \arrow[d, dashed]\\
										   &C
	\end{tikzcd}\]
	Note that this is the same as the diagram for cokernel, but rearranged slightly:
	\[\begin{tikzcd}[row sep = 2.5em, column sep = 2.5em]
		A_1\arrow[rr, bend left, "0"]\arrow[r, "\phi "']\arrow[dr, "0"']&A_2\arrow[r, "\psi "']\arrow[d, "c_2"]&C\\
									      &\varinjlim \mathscr{F} \arrow[ur, dashed]
	\end{tikzcd}\]
	Therefore, this definition of cokernel agrees with the preivously given definitions.
\end{proof}
\begin{exercise}
	Let $\phi : A \to B$ be a morphism in an additive category $\mathsf{A} $. Prove that $\iota : K \to A$ is a kernel for $\phi $ if and only if for all objects $Z$, the induced sequence
	\[\begin{tikzcd}[row sep = 2.5em, column sep = 2.5em]
		0\arrow[r]&\Hom_{\mathsf{A} }(Z, K) \arrow[r]&\Hom_{\mathsf{A} }(Z, A) \arrow[r]&\Hom_{\mathsf{A} }(Z, B) 
	\end{tikzcd}\]
	is exact. Formulate an analogous result for cokernels.
\end{exercise}
\begin{proof}
	The analogous result for cokernels is simply that
	\[\begin{tikzcd}[row sep = 2.5em, column sep = 2.5em]
		\Hom_{\mathsf{A} }(Z, B) \arrow[r]&\Hom_{\mathsf{A} }(Z, A) \arrow[r]&\Hom_{\mathsf{A} }(Z, K) \arrow[r]&0
	\end{tikzcd}\]
	is exact, which follows from our proof here as the dual of this statement.

	Let $\iota =\ker \phi $. Then we immediately have that since kernels are mono, $\zeta \circ \iota =0$ if and only if $\zeta =0$, and thus $\ker \iota =0$, and the complex is exact at $\Hom_{\mathsf{A} }(Z, B) $. It suffices to show $\im \iota =\ker \coker \iota =\ker \phi $. Let $\psi : \Hom_{\mathsf{A} }(Z, X)  \to \Hom_{\mathsf{A} }(Z, A) $ such that $\phi \circ \psi =0$. Then
	\[\begin{tikzcd}[row sep = 2.5em, column sep = 2.5em]
		Z\arrow[drr, dashed]\arrow[r, "\zeta "]\arrow[rrr, bend left, "0"]&X\arrow[r, "\psi "]&A\arrow[r, "\phi "]&B\\
							       &&K\arrow[u, "\iota "]
	\end{tikzcd}\]
	commutes, since $\ker \phi =\iota $. This means that the morphism $\iota : \Hom_{\mathsf{A} }(Z, K)  \to \Hom_{\mathsf{A} }(Z, A) $ is initial with respect to zero composition, and $\iota $ is $\ker \phi $. Since $\ker \phi $ exists, we have that $\ker \coker \ker \phi =\ker \phi $, and thus $\im \ker \phi =\ker \phi $, and the sequence is exact.

	The converse follows immediately from the diagram above.
\end{proof}
\begin{exercise}
	Let $\mathsf{A} $ be an additive category.
	\begin{itemize}
		\item Show that kernels are mono;
		\item Show that morphisms $\phi : A \to B$ with cokernels are epi if and only if $B\to 0$ is the cokernel;
		\item If $\mathsf{A} $ is abelian, prove that every kernel is the kernel of its cokernel.
	\end{itemize}
\end{exercise}
\begin{proof}
	I did these all in my notes instead of doing the book's arguments.
\end{proof}
\begin{exercise}
	Prove that for any abelian category $\mathsf{A} $, $\mathsf{A}^{op} $ is abelian.
\end{exercise}
\begin{proof}
	Finite products and coproducts exist because they are dual notions. The zero object remains zero, and monos become epis, which coencides with $\ker \coker $ becoming $\coker \ker $. The converse also holds. Abelian group structure holds since all that happens is arrows flip, and the structure remains the same. By the same logic, composition is still bilinear over addition.
\end{proof}
\begin{exercise}
	Let $\mathsf{A} $ be abelian and $\mathsf{C} $ be small. Show that $\mathsf{A} ^{\mathsf{C} } $ is abelian, and show that for any object $X$ in $\mathsf{C} $, the assignment $\mathscr{X} : \mathsf{A} ^{\mathsf{C} }  \to \mathsf{A} $ given by $\mathscr{F} \mapsto \mathscr{F} (A)$ is an exact functor.
\end{exercise}
\begin{proof}
	Recall that a natural transformation $\eta : \mathscr{F}  \Rightarrow \mathscr{G} $ is a family of morphisms $\eta_X : \mathscr{F} (X) \to \mathscr{G} (X)$ in $\mathsf{A} $, such that for any $\alpha : X \to Y$ in $\mathsf{C} $, the diagram
	\[\begin{tikzcd}[row sep = 2.5em, column sep = 2.5em]
		\mathscr{F} (X)\arrow[r, "\mathscr{F} (\alpha )"]\arrow[d, "\eta _X"']&\mathscr{F} (Y)\arrow[d, "\eta_Y"]\\
		\mathscr{G} (X)\arrow[r, "\mathscr{G} (\alpha )"']&\mathscr{G} (Y)
	\end{tikzcd}\]
	commutes.

	First, we show preadditivity. let $\nu ,\mu : \mathscr{F}  \Rightarrow \mathscr{G} $ be natural transformations. All properties other than closure under addition follow immediately from $\mathsf{A} $ being abelian, so it suffices to show $\nu+\mu $ is a natural transformation. Let $\alpha : X \to Y$ in $\mathsf{C} $. Since $\mathsf{A} $ is abelian and $\nu_X,\mu _X,\mathscr{F} (\alpha ),\mathscr{G} (\alpha )$ are morphisms in $\mathsf{A} $, by commutativity of the diagram since $\nu ,\mu $ are natural transformations, we have
	\[
		\left( \nu _X+\mu _X \right) \circ \mathscr{G} (\alpha )=\nu _X\circ \mathscr{G} (\alpha )+\mu _X\circ \mathscr{G} (\alpha )=\mathscr{F} (\alpha )\circ \nu _Y+\mathscr{F} (\alpha )\circ \mu _Y=\mathscr{F} (\alpha )\circ \left( \nu _Y+\mu _Y \right) 
	.\]
	Therefore, $\nu +\mu $ defined in this manner is a natural transformation $ \mathscr{F}  \Rightarrow \mathscr{G} $. Now let $\eta : \mathscr{G}  \Rightarrow \mathscr{H} $ be another natural transformation, and consider the composition $\eta \circ (\nu +\mu )$. Since $\nu +\mu $ is a natural transformation, so is $\eta \circ (\nu +\mu )$. Moreover, since $\eta _X,\nu _X,\mu _X$ are morphisms in $\mathsf{A} $, we have $\eta _X\circ (\nu _X+\mu _X)=\eta _X\circ \nu _X+\eta _X\circ \mu _X$. Therefore, $\mathsf{A} ^{\mathsf{C} } $ is preadditive.

	Now we need to show $\mathsf{A} ^{\mathsf{C} } $ is additive; that is, it has a zero object, finite products, and finite coproducts. Since $\mathsf{A} $ is abelian, it has a zero object, and thus there exists a 0 functor $\mathsf{C} \to \mathsf{A} $ given by $X\mapsto 0$, with all morphisms mapping to the unique morphism $0\to 0$. Note that any natural transformation $ \mathscr{F}  \Rightarrow 0$ or $ 0 \Rightarrow \mathscr{F} $ is manifestly unique, since any morphism $\mathscr{F} (X)\to 0$ or $0\to \mathscr{F} (X)$ in $\mathsf{A} $ is unique, as $0$ is initial and final. Therefore, $\mathsf{A} ^{\mathsf{C} } $ has a zero-object.

	Let $\mathscr{F} ,\mathscr{G} $ be functors $\mathsf{C} \to \mathsf{A} $, and let $\mathscr{F} \times \mathscr{G} : \mathsf{C}  \to \mathsf{A} $ be the functor $X\mapsto \mathscr{F} (X)\times \mathscr{G} (X)$. This manifestly satisfies the universal property for products, since there exist projections $\pi_{\mathscr{F} } ,\pi_{\mathscr{G} } : \mathscr{F} \times \mathscr{G}  \Rightarrow \mathscr{F} ,\mathscr{G} $ given by $\pi_{\mathscr{F}(X)} $ as the natural transformation $\mathscr{F} (X)\times \mathscr{G} (X)\to \mathscr{F} (X)$, and $\pi_{\mathscr{G} (X)} $ defined similarly. Then for any other natural transformation that gives morphisms $Z\to \mathscr{F} (X)$, there exists a unique morphism making the relevant diagram commute, since $\mathscr{F} (X)\times \mathscr{G} (X)$ is a product in $\mathsf{A} $. By similar logic, the functor $\mathscr{F} \amalg \mathscr{G} : \mathsf{C}  \to \mathsf{A} $ given by $X\mapsto \mathscr{F} (X)\amalg \mathscr{G} (X)$ is a coproduct. Therefore, $\mathsf{A} ^{\mathsf{C} } $ is additive.

	Finally, let $\eta : \mathscr{F}  \to \mathscr{G} $ be a natural transformation, and consider the morphism $\ker \eta : \mathscr{K}  \to \mathscr{F} $, defined by $\left( \ker \eta  \right)_X=\ker \left( \eta _X \right) $. It follows fairly immediately that the 0 natural transformation consists of 0 morphisms, so the diagram
	\[\begin{tikzcd}[row sep = 2.5em, column sep = 2.5em]
		\mathscr{Z} \arrow[r, Rightarrow]\arrow[rr, Rightarrow, bend left, "0"]\arrow[dr, Rightarrow, dashed]&\mathscr{F} \arrow[r, Rightarrow, "\eta "']&\mathscr{G} \\
													       &\mathscr{K} \arrow[u, Rightarrow, "\ker \eta "']
	\end{tikzcd}\]
	commutes if and only if
	\[\begin{tikzcd}[row sep = 2.5em, column sep = 2.5em]
		\mathscr{Z} (X)\arrow[rr, bend left, "0"]\arrow[dr, dashed]\arrow[r]&\mathscr{F}(X) \arrow[r, "\eta _X"']&\mathscr{G} (X)\\
										    &\mathscr{K} (X)\arrow[u, "\ker \eta _X"']
	\end{tikzcd}\]
	commutes for all $X$ in $\mathsf{A} $. This follows since kernels exist in $\mathsf{A} $, and thus $\ker \eta $ is the kernel of $\eta $. By the same logic, the obvious definition for $\coker \eta $ also satisfies the universal property for cokernels. The only difficult thing to show here is that the unique morphisms form a natural transformation, but this is obvious by commutativity of the diagrams.

	Now we have to show monomorphisms are kernels, and epimorphisms are cokernels. This is trivial to show (aka I don't want to type it up). Therefore, $\mathsf{A} ^{\mathsf{C} } $ is abelian.

	To show $\mathscr{X} $ is exact, let
	\[\begin{tikzcd}[row sep = 2.5em, column sep = 2.5em]
		0\arrow[r]&\mathscr{F} \arrow[r, "\mu "]&\mathscr{G} \arrow[r, "\nu "]&\mathscr{H} \arrow[r]&0
	\end{tikzcd}\]
	be exact. It follows that $\mu $ is mono and $\nu $ is epi, and thus $\im \mu =\ker \nu $ implies that $\ker \coker \mu =\mu =\ker \nu $, and similarly $\nu =\coker \mu $. It follows by our previous definition that $\mu_X=\ker \nu _X$, and $\coker \mu _X=\nu _X$, and thus
	\[\begin{tikzcd}[row sep = 2.5em, column sep = 2.5em]
		\mathscr{X}(0)\arrow[r]&\mathscr{X} (\mathscr{F} )\arrow[r, "\mathscr{X} (\mu )"]&\mathscr{X} (\mathscr{G} )\arrow[r, "\mathscr{X} (\nu )"]&\mathscr{X} (\mathscr{H} )\arrow[r]&0
	\end{tikzcd}\]
	is exact.
\end{proof}
\begin{exercise}
	For a ring $R$, let $\mathsf{R} $ denote the preadditive category with one element determined by $R$. Prove that $\Rmod{R} $ is equivalent to the full subcategory of \emph{additive} functors in $\mathsf{Ab} ^{\mathsf{R} } $. What about $\mathsf{Mod} $-$R$?
\end{exercise}
\begin{proof}
	a
\end{proof}
\begin{exercise}
	Let $R$ be a commutative ring, and let $\Rmod{R} ^{f} $ denote the category of finitely generated $R$-modules. Show that $\Rmod{R} $ need not be abelian, however if $R$ is Noetherian, then $\Rmod{R} $ is abelian.
\end{exercise}
\begin{proof}
	This follows immediately due to the following argument. There exists finitely generated $R$-modules $M$ over a general ring $R$ such that submodules of $M$ are not finitely generated. As such, kernels need not exist, since if the kernel of a morphism in $\Rmod{R} ^{f} $ need not be finitely generated.

	Conversely, if $R$ is Noetherian, then we have shown previously that all submodules of $M$ are Noetherian, and thus finitely generated. This concludes the proof.
\end{proof}
\begin{exercise}
	Let $T$ be a topological space, and recall that a presheaf on $T$ with values in a category $\mathsf{A} $ is a contravariant functor $\mathsf{T} \to \mathsf{A} $, where $\mathsf{T} $ is defined to be the category consisting of open sets in $T$, and morphisms corresponding to inclusions. That is, there exists $U\to V$ if $U\subseteq V$. Define the category of $A$-valued presheaves on $T$. Prove that if $\mathsf{A} $ is abelian, then the category of $\mathsf{A} $-valued presheaves on $T$ is abelian.
\end{exercise}
\begin{proof}
	This follows from exercise 9.1.1. Since a presheaf is a contravariant functor $\mathsf{T} \to \mathsf{A} $, the category of presheaves is simply $\mathsf{A} ^{\mathsf{T} ^{op} } $. By exercise 9.1.1, if $\mathsf{A} $ is abelian, then $\mathsf{A} ^{\mathsf{T} ^{op} } $ is abelian, since $\mathsf{T} $ is locally small.
\end{proof}
\begin{exercise}
	Consider the presheaf $\mathscr{C} $ of continuous complex-valued functions on a circle $S^1$, and the presheaf $\mathscr{C} ^{*} $ of continuous complex-valued functions that do not vanish anywhere along the circle. The first is a sheaf of abelian groups under $+$, and the second is a sheaf of abelian groups under $\cdot $. Use the exponential map to define a morphism $\exp : \mathscr{C}  \to \mathscr{C} ^*$, and show that $\exp $ is \emph{not} an epimorphism of presheaves.
\end{exercise}
\begin{proof}
	Recall the definition of these sheaves. For any open set $U$, we take $\mathscr{C} (U)$ to be the additive abelian group of functions $f : U \to \mathbb{C} $, and restructions $\mathscr{C} (V)\to \mathscr{C} (U)$ is taken to be the ordinary restriction of functions. We define $\mathscr{C} ^*$ similarly, and take our open sets to be the open sets on $S^1$ under the standard topology.

	A morphism of presheaves is simply a natural transformation $\exp : \mathscr{C}  \Rightarrow \mathscr{C} ^*$. We define this transformation in the natural way. Since $\mathscr{C}(U),\mathscr{C}^*\left( U \right) $ are abelian groups, we can define $\exp : \mathscr{C} (U) \to \mathscr{C} ^*(U)$ as the map $f\mapsto \exp \circ f$.

	I do not know how to show $\exp $ is not epi; it clearly uses the continuous hypothesis but I am not sure how so.
\end{proof}
\section{Working in Abelian Categories}
\section{Complexes and Homology, Again}
\begin{exercise}
	-
\end{exercise}
\begin{exercise}
	-
\end{exercise}
\begin{exercise}
	Verify that $\mathsf{C} (\mathsf{A} )$ is an abelian category.
\end{exercise}

\end{document}
